\section{At the Lord's Altar}

The Latin Church has long treasured a death told by a son. In the ninth book of the \emph{Confessions} Augustine recounts how, in the year of his baptism, he and his mother Monica waited at Ostia, at the mouth of the Tiber, for a ship to carry them home to Africa. ``As the day now approached on which she was to depart this life (which day Thou knewest, we did not), it fell out---Thou, as I believe, by Thy secret ways arranging it---that she and I stood alone, leaning in a certain window, from which the garden of the house we occupied at Ostia could be seen.''\endnote{Augustine, \emph{Confessions} IX.10.23, trans. J. G. Pilkington, \emph{Nicene and Post-Nicene Fathers}, first series, vol.~1.} The day was known to God and hidden from them. Neither of them was thinking of death. They were speaking of heaven.

``We were seeking between ourselves in the presence of the Truth, which Thou art, of what nature the eternal life of the saints would be, which eye hath not seen, nor ear heard, neither hath entered into the heart of man.''\endnote{\emph{Confessions} IX.10.23, Pilkington.} Their talk rose through every created thing, past the earth and the heavens and their own minds, toward the Wisdom by whom all things were made, ``and while we were thus speaking, and straining after her, we slightly touched her with the whole effort of our heart.''\endnote{\emph{Confessions} IX.10.24, Pilkington.} For a moment they knew what it is for the soul to be silent before God alone, and they asked whether that was not the meaning of the Lord's word, ``Enter thou into the joy of Thy Lord.'' Then they came back ``to the noise of our own mouth, where the word uttered has both beginning and end.''\endnote{\emph{Confessions} IX.10.24--25, Pilkington.}

Monica spoke first. ``Son, for myself, I have no longer any pleasure in aught in this life. What I want here further, and why I am here, I know not, now that my hopes in this world are satisfied. There was indeed one thing for which I wished to tarry a little in this life, and that was that I might see thee a Catholic Christian before I died. My God has exceeded this abundantly, so that I see thee despising all earthly felicity, made His servant,---what do I here?''\endnote{\emph{Confessions} IX.10.26, Pilkington.} Here is the whole of a good death in a few lines, spoken by a woman in good health. She had finished the work God gave her. Nothing held her in the world, and she was not afraid to leave it.

Within five days or not much more she fell ill of a fever. She swooned, came to herself, and looked at her two sons standing beside her. ``Here,'' she said, ``shall you bury your mother.'' Augustine's brother wished aloud that she might die in her own country. She silenced him with a look and spoke to them both: ``Lay this body anywhere, let not the care for it trouble you at all. This only I ask, that you will remember me at the Lord's altar, wherever you be.''\endnote{\emph{Confessions} IX.11.27, Pilkington.} She had once been anxious to lie beside her husband in Africa. Now the place of her grave did not matter to her, and when friends asked whether she did not dread leaving her body so far from her own city, she answered, ``Nothing is far to God; nor need I fear lest He should be ignorant at the end of the world of the place whence He is to raise me up.''\endnote{\emph{Confessions} IX.11.28, Pilkington.} On the ninth day of her sickness ``was that religious and devout soul set free from the body.''\endnote{\emph{Confessions} IX.11.28, Pilkington.}

Everything the Church teaches about Christian death is in that room. There is the soul's longing for God, already tasted at the window. There is the indifference to everything but God, which the masters of the spiritual life spend years trying to teach. There is the body entrusted without anxiety to the God who will raise it, and the one request a dying Christian makes of those who love her, which is to be remembered at the altar. The prayer of the Church for her dead, offered above all in the Sacrifice of the Mass, is expounded in \emph{What Death Does Not Finish}. Monica, at the end, wanted nothing else.

Augustine records his own grief as honestly as her death. He closed her eyes and held back his tears through the burial, because ``she neither died unhappy, nor did she altogether die.''\endnote{\emph{Confessions} IX.12.29, Pilkington.} Later, lying alone upon his bed, he let the memory of her return, ``and it was pleasant to me to weep in Thy sight, for her and for me, concerning her and concerning myself.'' He asks the reader not to laugh at him for weeping ``for my mother during so small a part of an hour,---that mother who was for a while dead to mine eyes, who had for many years wept for me, that I might live in Thine eyes.''\endnote{\emph{Confessions} IX.12.33, Pilkington.} Ambrose's lesson holds. The son who believed she had not altogether died still wept for her, and his tears were a prayer.
